\chapter{Épreuve chronométrée 3}\label{ch:iv1:timed-test-3}

\begin{ivtrains}
Il faut la traiter d'une seule traite, sans notes ni calculatrice.
\end{ivtrains}

\ivtestsection{A}{6}{40}

\begin{question}\label{iq:iv1:timed-test-3:1}
Une boîte de céréales contient l'un de trois jouets, avec les probabilités $\frac12$, $\frac13$ et $\frac16$, indépendamment d'une boîte à l'autre. En moyenne, combien de boîtes achetez-vous pour obtenir les trois~?
\end{question}

\begin{question}\label{iq:iv1:timed-test-3:2}
On lance trois dés équilibrés. Lequel est le plus probable, un total de 9 ou un total de 10~? Donnez les deux probabilités.
\end{question}

\begin{question}\label{iq:iv1:timed-test-3:3}
Une source émet quatre symboles avec les probabilités $\frac12$, $\frac14$, $\frac18$ et $\frac18$. Quel est le plus petit nombre moyen de bits par symbole que peut atteindre un code préfixe~? Quelle est l'entropie de la source~?
\end{question}

\begin{question}\label{iq:iv1:timed-test-3:4}
$X$ suit une loi uniforme sur $[0, 1]$, et sachant $X$, $Y$ suit une loi uniforme sur $[0, X]$. Quelle est la probabilité que $Y > \frac12$~?
\end{question}

\begin{question}\label{iq:iv1:timed-test-3:5}
Deux joueurs lancent un dé à tour de rôle, A commençant. A gagne en obtenant un 6. B gagne en obtenant un 5 ou un 6. Quelle est la probabilité que A gagne~?
\end{question}

\begin{question}\label{iq:iv1:timed-test-3:6}
On choisit au hasard trois sommets d'un hexagone régulier. Quelle est la probabilité qu'ils forment un triangle rectangle~? Et un triangle équilatéral~?
\end{question}

\ivtestsection{B}{6}{60}

\begin{question}\label{iq:iv1:timed-test-3:7}
On choisit au hasard une permutation de $1, \dots, 8$. Quelle est la probabilité que l'élément 8 appartienne à un cycle de longueur exactement 3~? Quel est le nombre moyen de cycles~?
\end{question}

\begin{question}\label{iq:iv1:timed-test-3:8}
Une maladie touche 1~\% des personnes, indépendamment les unes des autres. Vous regroupez le sang de 10 personnes et testez le mélange une seule fois. Si le résultat est positif, vous testez les 10 personnes individuellement. Quel est le nombre moyen de tests par personne~?
\end{question}

\begin{question}\label{iq:iv1:timed-test-3:9}
Vous lancez sans cesse un dé équilibré et tenez à jour le total cumulé. Quelle est la probabilité que le total cumulé soit un jour égal à exactement 6~? Quelle est, à long terme, la chance d'atteindre un grand nombre $n$~?
\end{question}

\begin{question}\label{iq:iv1:timed-test-3:10}
Un bâton de longueur 1 est cassé en un point choisi uniformément au hasard, et les deux morceaux deviennent les côtés d'un rectangle. Quelle est l'aire moyenne du rectangle~? Quelle est la probabilité que l'aire dépasse $\frac{3}{16}$~?
\end{question}

\begin{question}\label{iq:iv1:timed-test-3:11}
Un call à la monnaie de trois mois sur une action cotée 100, avec une volatilité de 16~\%, vaut environ $0.4S\sigma\sqrt{T}$, les taux étant nuls. Estimez le prix du straddle. Quel mouvement quotidien permet à un détenteur couvert en delta d'atteindre le seuil de rentabilité, en prenant 252 jours de cotation par an~?
\end{question}

\begin{question}\label{iq:iv1:timed-test-3:12}
Une permutation aléatoire de $1, \dots, 10$ est écrite en ligne. Une position est un maximum local si sa valeur dépasse celles de ses voisins, chaque extrémité n'ayant qu'un seul voisin. Combien y a-t-il de maxima locaux en moyenne~?
\end{question}
